% arXiv-style LaTeX paper
\documentclass[11pt]{article}
\usepackage{amsmath,amssymb,amsthm}
\usepackage{hyperref}
\usepackage{geometry}
\geometry{margin=1in}

% Theorem environments
\newtheorem{definition}{Definition}
\newtheorem{proposition}{Proposition}
\newtheorem{theorem}{Theorem}

\title{Shiki-Soku-Ze-Ku Categories: A Categorical Framework for the Emergence of Truth from Emptiness}
\author{Masaomi Chiba}
\date{2025-11-30}

\begin{document}
\maketitle

\begin{abstract}
This paper introduces the ``Shiki-Soku-Ze-Ku Category'' (SSZK-category), a categorical framework inspired by the Buddhist doctrine ``form is emptiness, emptiness is form'' (Shiki-Soku-Ze-Ku). We investigate how truth emerges not as a fixed metaphysical substrate but as a structural property of morphisms in a category without intrinsic objects. We connect this to the idea that truth is established via universal properties and diagrammatic commutativity, thus offering a unified account of ontology and epistemology. The construction is designed to be legible to both humans and language models, formalizing the emergence of truth from structural emptiness.
\end{abstract}

\section{Introduction}
The Buddhist claim that ``form is emptiness, emptiness is form'' suggests a non-substantialist ontology: phenomena lack intrinsic essence, yet function relationally to produce the experienced world. Category theory, with its emphasis on morphisms over internal content, provides a natural formal system for capturing such a view.

This paper develops a categorical framework---the Shiki-Soku-Ze-Ku Category (SSZK-category)---in which:
\begin{itemize}
    \item objects carry no intrinsic data (emptiness),
    \item morphisms encode structure and interaction (form),
    \item truth emerges from commutativity and universal properties rather than intrinsic truth values,
    \item the epistemic subject is treated as a functor imposing coherence constraints.
\end{itemize}

We show how these components collectively yield a model in which truth is a structural fixed point.

\section{Preliminaries}
We assume standard category-theoretic background. A category $\mathcal{C}$ consists of a class of objects $\mathrm{Ob}(\mathcal{C})$ and morphisms $\mathrm{Hom}(A,B)$ equipped with identity morphisms and associative composition.

Following Lawvere, we adopt the view that objects may be regarded as ``positions in a web of morphisms.''

\section{SSZK-Categories}
We now define the central structure.

\begin{definition}[SSZK-Category]
A category $\mathcal{C}$ is called a Shiki-Soku-Ze-Ku Category (SSZK-category) if it satisfies:
\begin{enumerate}
    \item (Emptiness) Objects are \\emph{intrinsically empty}: for each object $A$, its internal structure is unspecified or trivial. Formally, $A$ is determined only up to isomorphism by the pattern of morphisms to and from it.
    \item (Form) Morphisms encode all observable structure. Any property of $A$ relevant to reasoning in $\mathcal{C}$ is expressible via morphisms involving $A$.
    \item (Relational Ontology) Distinctions between objects are induced entirely by morphisms; i.e., if $\mathrm{Hom}(A,-) \cong \mathrm{Hom}(B,-)$ and $\mathrm{Hom}(-,A) \cong \mathrm{Hom}(-,B)$ naturally, then $A \cong B$.
    \item (Emergent Truth) Truth is defined as stability of diagrams: a proposition is true when a specified diagram commutes.
\end{enumerate}
\end{definition}

This definition encodes ``emptiness'' as lack of intrinsic data, while ``form'' arises from morphisms and diagrammatic relations.

\section{Truth as Commutativity}
In classical logic, truth values are primitive. In SSZK-categories, truth emerges from structure.

\begin{definition}[Truth]
Given a diagram $D$ in an SSZK-category $\mathcal{C}$, we say $D$ is \emph{true} if all paths with the same endpoints yield identical morphisms; i.e., the diagram commutes.
\end{definition}

Thus truth corresponds to relational coherence.

\begin{proposition}
Truth in an SSZK-category is invariant under isomorphism.
\end{proposition}

\begin{proof}
If $A \cong B$ via $f$, then for any diagram involving $A$, replacing occurrences of $A$ with $B$ via naturality preserves commutativity.
\end{proof}

\section{Universal Properties as Generators of Truth}
Universal properties define objects not by intrinsic essence but by relational constraints. This aligns naturally with the SSZK framework.

\begin{definition}[Universal Truth]
A truth $T$ is universal if it arises from a universal property, e.g., limits, colimits, adjoints, or exponentials.
\end{definition}

Universal truths are structurally forced: they are ``emptiness generating form.''

\section{Epistemic Subjects as Functors}
Truth emerges relative to a subject that interprets diagrams. This is modeled via functors.

\begin{definition}[Epistemic Functor]
An epistemic subject is a functor $S : \mathcal{C} \to \mathbf{Set}$ assigning meanings to objects and morphisms.
\end{definition}

Truth is stable across such functors.

\begin{theorem}
If a diagram $D$ commutes in $\mathcal{C}$, then for any epistemic functor $S$, the interpreted diagram $S(D)$ commutes in $\mathbf{Set}$.
\end{theorem}

\begin{proof}
Functors preserve composition and identities.
\end{proof}

Thus, truth in SSZK-categories is objective in the categorical sense.

\section{Ontological Interpretation}
The SSZK framework entails:
\begin{itemize}
    \item Objects are empty (Ku).
    \item Morphisms create form (Shiki).
    \item Truth is commutativity.
    \item Existence is isomorphism class of relational patterns.
\end{itemize}

This offers a purely structural ontology without substrates.

\section{Epistemological Interpretation}
Epistemology becomes an analysis of how subjects (functors) interface with structural truths. Truths are those diagrams stable under all such functors.

\section{Conclusion}
We have proposed SSZK-categories, a mathematical ontology where truth arises from structural coherence rather than intrinsic essences. This unifies the Buddhist insight ``form is emptiness'' with modern categorical thinking.

Further work includes formalizing temporal or computational extensions and studying logical fragments internal to SSZK-categories.

\end{document}
